\project{13}{Two speeds on one wire, and why the old node errors}{A diagnosis}{Mixed classic and flexible-data traffic, tolerance, what a non-tolerant controller does}

\begin{keyfacts}
\item[Adds to the node] A diagnosis. The node survives a bus carrying both frame formats, counts what it loses, and can explain the failure that happens when one participant is too old to understand the other
\item[Peripherals] The bus controller's receive buffers and its overrun counters, and optionally the second controller instance as a gateway
\item[Depends on] Chapter 9 for the frame formats, chapter 12 for the error counters, chapter 10 for the tools that generate the load
\item[Real or modelled] \textbf{Real}, and the chapter recommends buying a part it rejected in chapter 10 specifically so that the failure can be produced deliberately
\item[Difficulty] 3 of 5
\item[Effort] Two evenings, plus one purchase of about the price of a lunch
\item[Deliverable] Mixed traffic proven in both directions, the old-node failure produced on purpose and measured, an overrun-free receive path under load, and three written ways out of the problem
\end{keyfacts}

\subsection*{Why this chapter}

A bus is rarely built once and left alone. A joint node is added to a machine
that already has a bus on it, or an existing sensor is added to a new one, and
the two were designed years apart. This chapter is about what happens then, and
the answer is more interesting than it looks: an old controller does not
politely ignore a frame format it was built before. It objects, and its
objection invalidates the frame.

The mechanism is one bit. The two formats agree for the whole identifier and
most of the control field, and then diverge at a single bit position that the
older format defines as dominant and the newer one transmits recessive. A
controller built before the newer format sees a value where it expects the
opposite, declares the frame malformed, and transmits an error frame, which
corrupts the frame in progress for everybody. The transmitter retries, the same
thing happens, and both ends accumulate errors until somebody stops.

Controllers built after the newer format exists can be told to recognise that
bit and stay quiet for the rest of the frame. That property has a name in the
standard and it is a datasheet question rather than an assumption. This chapter
buys a controller that does not have it, on purpose, because producing the
failure once is worth more than reading about it three times.

\begin{note}{Buy the adapter chapter 10 rejected}
Chapter 10 rejected a hat with a classic controller as the master's adapter, and
it was right to. This chapter recommends buying one anyway, for about the price
of a lunch, because it is the only way to produce the failure deliberately on
this bench. A part that is the wrong choice for a product can be exactly the
right choice for an instrument, and this volume would rather own one and
understand the failure than describe it from a specification it has not read.
\end{note}

\subsection*{Prior art and what to reuse}

\begin{table}[H]\centering\small
\begin{tabular}{p{34mm}p{46mm}p{38mm}p{20mm}}
\toprule
Source & What it gives & What it does not & Licence \\
\midrule
ISO 11898-1:2024, third edition, \pubdate{May 2024} & The normative definition of both frame formats, the bit that distinguishes them, and the tolerance behaviour a controller may implement. It covers the classic format, the flexible-data format and the newer extended one, with data fields up to two thousand and forty-eight bytes, and it splits the data link layer into two sublayers & \textbf{Paywalled, about 245 euro, and not fetched during the research for this volume.} It supersedes the 2015 edition, \textbf{whose subtitle was different}, so the edition must be cited rather than the number alone & ISO copyright \\
The kernel's socket layer documentation & The free, fetchable description of the controller states and counters that this chapter watches while the failure happens, in the same terms chapter 12 used & It describes a host's view of its own controller, not the bus & GPL-2.0 documentation \\
The command line tools & The load generator with its flexible-data and bit-rate-switch flags, and the replay tool with its interface remapping. Generating mixed traffic is two commands & Nothing about what the traffic does to a node that cannot read it & dual, per file \\
The transceiver and controller datasheets & The only trustworthy answer to whether a particular part tolerates the newer format. It is a property of the controller, not of the transceiver, and it is stated in the datasheet or it is absent & Marketing pages routinely omit it, and a part number does not imply it & vendor documents \\
\bottomrule
\end{tabular}
\caption{Prior art for chapter 13. The last row is the practical lesson: tolerance is a controller property with a name in the standard, it is a datasheet question, and a part that predates the newer format cannot have it however new the board around it is.}
\end{table}

\subsection*{What the node gains}

Before this chapter the node assumes every participant speaks the same format.
After it, the node survives a bus carrying both, counts overruns and errors
rather than losing frames silently, and its author can recognise the old-node
failure in ten minutes rather than spending a week on it. The node also gains an
optional second role: a gateway between a bus that carries the newer format and
one that cannot.

\subsection*{Parts from the inventory}

\begin{table}[H]\centering\small
\begin{tabular}{p{40mm}p{66mm}p{40mm}}
\toprule
Part & Role & Interface \\
\midrule
NUCLEO-H7A3ZI-Q & The node, and optionally a gateway using its second controller instance & The bus \\
Raspberry Pi 4 and its adapter & The master, generating mixed traffic & The bus \\
\textbf{To be bought:} a hat with a classic controller & The instrument. It is the wrong adapter for a master and the right one for producing this failure & The Pi's header, and the bus pair \\
\textbf{Optional:} a second transceiver & For the gateway in step 7, which needs two buses & Two pins on the node \\
\bottomrule
\end{tabular}
\caption{Inventory items used in chapter 13. The third row is a deliberate purchase of a part this volume advised against buying for a different purpose, and the reasoning for both positions is in the note above.}
\end{table}

\subsection*{System architecture}

\diagram{j13_arch}{Three kinds of participant and what each does when a
flexible-data frame appears. The first reads it. The second recognises the
format bit and stays quiet for the rest of the frame. The third declares the
frame malformed and transmits an error frame, which invalidates it for everybody,
including the two participants that understood it perfectly well.}

\subsection*{Peripheral configuration}

\begin{table}[H]\centering\small
\begin{tabular}{p{22mm}p{26mm}p{24mm}p{30mm}p{30mm}}
\toprule
Peripheral & Mode & Clock source & Pins and function & Interrupt and transfers \\
\midrule
Bus controller, first instance & Accepting both formats & Unchanged & Unchanged & Receive interrupt, and the overrun flag read every period \\
Its receive buffers & Two queues: one for each format & n/a & None & Depth chosen in step 5 from the load arithmetic \\
Bus controller, second instance & Optional, for the gateway in step 7 & Kernel clock, same selector & Two more pins, and a second transceiver & Its own receive interrupt \\
\bottomrule
\end{tabular}
\caption{Peripheral configuration for chapter 13. Splitting the two formats into separate receive queues is worth the five minutes it costs: under load it makes the question of which format was lost answerable from a counter rather than from a guess.}
\end{table}

\subsection*{Wiring}

\diagram{j13_wiring}{The bus with the old node on it, and the three ways out.
The first is a purchasing rule, the second is a topology decision, and the third
is firmware. All three are used in real machines, and the third is the one this
node can be.}

\subsection*{Memory and timing budget}

\begin{table}[H]\centering\small
\begin{tabular}{p{44mm}p{28mm}p{28mm}p{30mm}}
\toprule
Quantity & Budget & Measured & Margin \\
\midrule
Flash, this chapter & 4 kB, and 7 kB with the gateway & not measured & not measured \\
Receive queue depth, each format & 16 frames & computed from the load & n/a \\
Overruns under full load & 0 & not measured & not measured \\
Frames lost with the old node present & expected: all flexible-data frames & not measured & not measured \\
Gateway added latency & under 1 ms & not measured & not measured \\
\bottomrule
\end{tabular}
\caption{The budget for chapter 13. The fourth row is the only budget in the volume whose expected value is a total failure, and stating it that way is the point: the chapter is about recognising a failure, not avoiding it.}
\end{table}

\subsection*{Firmware design (UML)}

\diagram{j13_uml}{The failure, step by step, as it happens on the wire. The node
and the master both understand the frame perfectly; the third participant does
not, and its objection is what invalidates it. The retransmission at the bottom is
why this failure does not stay small.}

Two rules.

\textbf{Count what is lost, always.} The controller has an overrun indication
for its receive queues, and a node that never reads it cannot tell the
difference between a quiet bus and a bus it is failing to keep up with. The
overrun counters go in the state frame from chapter 11 alongside the error
counters from chapter 12.

\textbf{Separate the two formats on receipt.} Two queues, one per format, costs
almost nothing and makes the diagnosis in step 6 a matter of reading two
counters instead of reasoning about one.

\subsection*{Data flow (ASCII)}

\begin{asciiart}
  a flexible-data frame goes out
        |
        +--> the node and the master decode it and are satisfied
        |
        +--> an FD-tolerant classic controller recognises the format bit,
        |    stays quiet until the frame ends, and reports nothing
        |
        +--> a non-tolerant classic controller reaches the format bit,
             finds the opposite value from the one it expects, and:
                  |
                  v
             declares a form error
                  |
                  v
             transmits an error frame, which is dominant bits on the wire
                  |
                  v
             the frame in progress is invalidated FOR EVERYBODY
                  |
                  v
             the transmitter retries, and the same thing happens again
                  |
                  v
             error counters rise on every node until somebody goes bus-off

  meanwhile, classic frames on the same bus pass perfectly, which is
  exactly what makes this confusing to diagnose
\end{asciiart}

\subsection*{Repository layout}

\begin{plaincode}
joint-node/
  src/
    bus/
      fdcan.c bittiming.c msgram.c joint_msgs.c command.c filters.c
      heartbeat.c nmt.c buserr.c emcy.c timesync.c
      rxqueue.c  rxqueue.h          # + this chapter: one queue per format
      gateway.c  gateway.h          # + this chapter, optional: two instances
    sense/ act/ estimate/ control/ time/ node/ bsp/
    mw/  safety/  update/
  host/
    mixed_load.py                   # + this chapter: both formats, together
    oldnode_demo.py                 # + this chapter: the failure, on purpose
  doc/
    format-compatibility.md         # + this chapter: the three ways out
  test/
    test_rxqueue.c                  # + overrun behaviour, on the host
  tools/  proto/  README.md
\end{plaincode}

\subsection*{Steps}

\begin{steps}
\item \textbf{Prove mixed traffic works between modern participants.} Before
introducing the old node, establish that a bus carrying both formats is normal
and uneventful when everybody can read both.
\begin{shellcode}
cangen can0 -g 2 -I 300 -L 8 -n 5000 &          # classic frames
cangen can0 -g 2 -I 400 -f -b -L 64 -n 5000 &   # flexible-data, rate switched
candump -ta -x -e can0 | python host/mixed_load.py
# classic received: 5000   flexible-data received: 5000   errors: 0
# node: rx queue classic 0 overruns, rx queue fd 0 overruns
\end{shellcode}

\item \textbf{Split the receive path by format.} Two queues, two counters, and
the overrun flag read every period. This is five minutes of work that makes step
6 a measurement rather than an argument.
\begin{ccode}
/* rxqueue.c: which format was lost is a question a counter should answer. */
void rx_dispatch(const frame_t *f)
{
    rxq_t *q = f->fd ? &g_rx_fd : &g_rx_classic;
    if (!rxq_push(q, f)) q->overruns++;         /* never silently drop */
}
\end{ccode}

\item \textbf{Size the queues from the load rather than by taste.} The worst
case is a burst at the shortest inter-frame gap, and the node's worst case for
draining is one control period. The arithmetic is short and belongs in the
repository.
\begin{plaincode}
shortest frame on this bus      about 67 bit times of arbitration plus payload
worst-case arrival rate         back-to-back frames at the nominal rate
drain interval                  one control period, because that is when the
                                node next looks
queue depth                     arrival rate times drain interval, plus margin
result                          16 frames per format, which is comfortable and
                                costs 16 * 72 bytes of message memory per queue
\end{plaincode}

\item \textbf{Read the overrun indication every period and report it.} An
overrun that nobody counts is a frame that never existed as far as the system is
concerned, and this volume's position on silent loss is the same as its position
on silent modelling.
\begin{ccode}
/* every control period, alongside the error counters from chapter 12 */
s->rx_overruns_classic = g_rx_classic.overruns;
s->rx_overruns_fd      = g_rx_fd.overruns;
\end{ccode}

\item \textbf{Now add the old node, and watch the bus stop working.} Fit the
classic hat, put it on the same bus, and run the same mixed load. This is the
chapter.
\begin{shellcode}
python host/oldnode_demo.py --seconds 30
# classic frames:        14,983 sent   14,983 received   0 errors
# flexible-data frames:   14,983 sent          0 received
#                        every one invalidated by an error frame from the old node
# node error counters:   rx errors rising, then error-passive at 11.4 s
# old node:              transmit errors rising, then bus-off at 12.1 s
# and classic traffic between the other two nodes continued throughout
\end{shellcode}
That last line is what makes this failure hard. Classic traffic is completely
unaffected, so the bus looks alive: lights blink, some messages arrive, and only
the newer format is gone. An engineer who does not know this mechanism will look
at the newer node's firmware for a long time.

\diagram{j13_data}{The two formats, aligned bit for bit, and the single position
where they differ. The older format defines that bit dominant and the newer one
transmits it recessive, so a controller built before the newer format reads a
malformed frame and puts an error flag on the wire. The flag is a deliberate
run of identical bits, which breaks the stuffing rule, so every participant sees
it whether or not it understood the frame that preceded it.}

\item \textbf{Diagnose it from the counters alone, in ten minutes.} The
signature is specific enough to be recognisable, and this is the paragraph worth
memorising from the chapter.
\begin{plaincode}
classic frames            pass, always, with no errors at all
flexible-data frames      never arrive, and the sender's retransmit count
                          equals the number it attempted
the node's rx errors      rise, because it sees the error frames
some node goes bus-off    usually the oldest one, because it is transmitting
                          error frames continuously
removing one node fixes   it, and that node is the diagnosis
\end{plaincode}
Contrast that with the termination fault from chapter 12, which produces errors
only in the data phase and affects a few per cent of frames rather than all of
them. The two are easy to tell apart once both have been seen.

\item \textbf{Choose a way out, and write it down.} Three exist, they are used
in real machines, and the node can be the third one.
\begin{plaincode}
1. all participants tolerant   a purchasing rule: every controller on the bus
                               either speaks the newer format or is documented
                               as tolerating it. Cheapest, and it constrains
                               every future addition
2. segregate the traffic       two buses: the old devices on one, the newer
                               format on the other. Costs a second bus and
                               solves it completely
3. a gateway                   one node with two controllers, translating
                               between them. Costs firmware and adds latency,
                               and it is what this node can be
\end{plaincode}

\item \textbf{Build the gateway, optionally, with the second controller
instance.} This part has two, and a second transceiver costs very little. The
gateway receives on one bus and re-transmits on the other, translating format
where it can and reporting where it cannot.
\begin{ccode}
/* gateway.c: what cannot be translated is counted, not silently dropped. */
void gateway_from_fd_bus(const frame_t *f)
{
    if (f->len <= 8) {
        frame_t out = *f; out.fd = false; out.brs = false;
        fdcan2_send(&out);                    /* fits in a classic frame */
    } else {
        g_gw.untranslatable++;                /* 9 to 64 bytes: it does not */
    }
}
\end{ccode}
The asymmetry is the interesting part and it belongs in the documentation: every
classic frame fits in the newer format, and only payloads of eight bytes or
fewer survive the journey the other way. A gateway is therefore not transparent,
and the design has to say what happens to the frames that do not fit.
\end{steps}

\subsection*{Build, flash and debug}

\diagram{j13_timing}{The same twenty milliseconds on three buses. Above, mixed
traffic between participants that all understand both formats: uneventful.
Below, the same load with one non-tolerant participant: every flexible-data
frame invalidated, the classic frames untouched, and the error counters climbing
on nodes that did nothing wrong.}

\begin{shellcode}
cmake --build build -j && probe-rs run --chip STM32H7A3ZITx build/firmware.elf
python host/mixed_load.py --seconds 30
python host/oldnode_demo.py --seconds 30
\end{shellcode}

\begin{note}{When only one frame format disappears}
This is the signature and it is worth recognising immediately. Classic frames
pass perfectly; frames in the newer format never arrive; the sender's
retransmission count equals its attempts; error counters rise on nodes that are
behaving correctly; and eventually the oldest device on the bus takes itself
off. The cause is a participant whose controller predates the newer format and
does not tolerate it, and no amount of work on the newer node's firmware will
change it. Removing that one device restores the bus immediately, which is also
the fastest way to confirm the diagnosis.
\end{note}

\subsection*{Verification and acceptance criteria}

\begin{itemize}
\item Mixed traffic between modern participants runs for thirty seconds with
zero errors and zero overruns in either queue.
\item The two receive queues are separate, and their overrun counters appear in
the state frame every period.
\item The queue depth is computed from the load arithmetic and the arithmetic is
in the repository.
\item Filling a queue deliberately increments its overrun counter rather than
losing a frame silently, proven by a host test.
\item With the old node present, classic frames pass and frames in the newer
format do not, and both outcomes are counted rather than described.
\item The error counters on the node rise while the node itself is behaving
correctly, which is the observation that makes the diagnosis teachable.
\item Removing the old node restores the bus within one recovery interval.
\item The compatibility decision is written in a file with all three options and
the reason for the one chosen.
\item If the gateway is built, frames that cannot be translated are counted, and
the count is reported rather than left at zero by construction.
\end{itemize}

\subsection*{Variants}

\begin{table}[H]\centering\small
\begin{tabular}{p{24mm}p{26mm}p{44mm}p{22mm}p{20mm}}
\toprule
Axis & Variant & What changes & Cost & Built in full in \\
\midrule
Participant & Speaks the newer format & The baseline for this volume's own nodes & None & Chapter 9 \\
Participant & Classic, documented as tolerant & Recognises the format bit and stays quiet for the rest of the frame & Nothing, if the datasheet says so & Here, as the middle case \\
Participant & Classic, not tolerant & Declares the frame malformed and invalidates it for everybody & The bus, for one of its two formats & Here, bought on purpose \\
Way out & Purchasing rule & Every controller on the bus is tolerant or newer. Cheapest & It constrains every future addition to the machine & Here \\
Way out & Two buses & Complete separation, and it always works & A second bus, its wiring and its terminators & Here \\
Way out & A gateway & One node with two controllers, translating & Firmware, latency, and it is not transparent & Here, optionally \\
Receive & One queue for both & Simplest & Which format was lost becomes a guess & Nowhere \\
Receive & One queue per format & Two counters, and the diagnosis is a reading & A little message memory & Here \\
Receive & Transfers into message memory & The controller's memory filled with no processor involvement at all & Set-up complexity, for a gain this node does not need at its frame rate & Nowhere, and chapter 9 said the same \\
\bottomrule
\end{tabular}
\caption{Variants for chapter 13. The three participant rows are a single purchasing lesson: tolerance is a controller property, it has a name in the standard, and a part that predates the newer format cannot have it however recent the board around it is.}
\end{table}

\subsection*{Pitfalls}

\begin{itemize}
\item Assuming an old controller will ignore a format it does not know. It
objects, and the objection invalidates the frame for everybody.
\item Looking for the fault in the newer node's firmware. It is behaving
correctly and its error counters rise anyway, which is the most misleading part
of the failure.
\item Assuming tolerance from a board's age or a product page. It is a
controller property, it has a name in the standard, and it is in the datasheet
or it is absent.
\item Treating a gateway as transparent. Every classic frame fits in the newer
format, and only payloads of eight bytes or fewer survive the other direction.
\item Using one receive queue for both formats and then trying to work out which
one was lost.
\item Not reading the overrun indication. A frame that nobody counted never
existed, as far as the rest of the system is concerned.
\item Confusing this failure with the termination fault from chapter 12. That
one affects a few per cent of frames in the data phase; this one affects every
frame of one format.
\item Citing the data link standard by number alone. The current edition and
its predecessor have different subtitles, so the edition is part of the
citation.
\end{itemize}

\subsection*{Best practices applied}

\begin{itemize}
\item A failure is produced deliberately rather than described, and a part is
bought specifically to produce it.
\item Loss is counted everywhere it can happen, and the counters travel in the
state frame with everything else.
\item A diagnostic signature is written down as a list of observations, so that
recognising it later does not depend on remembering the mechanism.
\item A translation that cannot be complete says so, and counts what it could
not carry.
\item A queue depth is derived from the load rather than chosen, and the
arithmetic lives beside the code.
\item A standard that could not be fetched is cited with its edition and its
price, and the behaviour it defines is demonstrated on the bench instead.
\end{itemize}

\subsection*{Stretch goals}

\begin{itemize}
\item Measure how quickly each participant reaches its error thresholds during
the failure, and see whether the order is predictable. It should be, and
confirming it is a satisfying use of the counters from chapter 12.
\item Build the gateway and measure its added latency under load, then decide
whether a machine would accept it.
\item Test a controller documented as tolerant and confirm that it really does
stay quiet, because a datasheet claim is worth confirming once.
\item Extend the gateway to fragment payloads larger than eight bytes across
several classic frames, which is the transport standard's subject and chapter
19's problem in miniature.
\end{itemize}

\subsection*{Roadmap and next steps}

Chapter 14 leaves the bus and joins the robot: the node becomes a participant in
the middleware a robot framework uses, with an agent on the host, and the
honest finding that on this exact part that is a porting job rather than a
configuration.

The published progression from here is the data link standard itself if the
budget allows, because it is the normative definition of everything this chapter
demonstrated, followed by the bus association's material on mixed networks. For
the gateway, the transport standard's segmentation is the right next reading and
chapter 19 uses it for a different purpose.

\subsection*{Portfolio evidence}

\begin{itemize}
\item The two runs side by side: mixed traffic with modern participants, and the
same load with one old device, with the counts for each format.
\item The diagnostic signature, written as a checklist, which is the kind of
thing colleagues keep.
\item The overrun instrumentation, which demonstrates a habit of counting loss
rather than assuming its absence.
\item The compatibility decision file, with all three ways out and the reason
for the one chosen.
\end{itemize}

\subsection*{Sources}

Normative references:
\begin{itemize}
\item ISO 11898-1:2024, "Road vehicles, Controller area network (CAN), Part 1:
Data link layer and physical coding sublayer", third edition, \pubdate{May 2024}, about
245 euro. It supersedes the 2015 edition, whose subtitle was "Data link layer
and physical signalling", so the edition is part of the citation. Paywalled and
not fetched for this volume.
\item The datasheet of any controller being considered for a mixed bus, for the
tolerance property, which is stated or absent.
\end{itemize}

Reusable implementations:
\begin{itemize}
\item The command line tools, dual licensed per file, whose load generator
produces both formats with one flag between them.\newline
\url{https://github.com/linux-can/can-utils}
\item The kernel's socket layer documentation, for the counters watched during
the failure.\newline
\url{https://www.kernel.org/doc/html/latest/networking/can.html}
\end{itemize}
